\documentclass[letterpaper,10pt]{article}

% Compiled with pdflatex (this app's backend always invokes pdflatex, not
% xelatex/lualatex) — so fontspec + system font loading was replaced below
% with pdflatex-compatible helvet, keeping the same clean sans-serif look.
\usepackage[utf8]{inputenc}
\usepackage[T1]{fontenc}
\usepackage{geometry}
\usepackage[hidelinks]{hyperref}
\usepackage{titlesec}
\usepackage{enumitem}
\usepackage{xcolor}
\usepackage{tabularx}
\usepackage{helvet}
\renewcommand{\familydefault}{\sfdefault}
% paracol (not minipage) for the two-column body: a minipage is a single
% atomic, unbreakable box, so once the combined column content is taller
% than the space left on the page after the header, LaTeX pushes the
% WHOLE block to the next page — leaving the header stranded alone on
% page 1. paracol's columns break across pages like ordinary text instead.
\usepackage{paracol}

% Precise Deedy Geometric Margins
\geometry{left=1.25cm, top=1cm, right=1.25cm, bottom=1cm}

% Without this, the narrow 0.33\textwidth left column gets fully
% justified by default, which stretches inter-word spacing unevenly on
% short lines (most visible in the Skills entries) — raggedright avoids
% that stretching everywhere, matching every other template in this set.
\raggedright
\raggedbottom

% Exact Deedy Color Palette Matrix
\definecolor{primarytext}{RGB}{45, 55, 72}
\definecolor{headings}{RGB}{107, 70, 193}     % Classic Deedy Purple Accent
\definecolor{subheadings}{RGB}{45, 55, 72}    % Deep Gray Titles
\definecolor{graytext}{RGB}{113, 128, 150}    % Soft Gray Metadata

% Header Styling Rules
\titleformat{\section}{\color{headings}\scshape\Large\bfseries}{}{0em}{}
\titlespacing*{\section}{0pt}{5pt}{2pt}

% Structure Component Directives
\newcommand{\descript}[1]{{\color{subheadings}\textit{\small #1}}}
\newcommand{\location}[1]{{\color{graytext}\footnotesize #1}}
\newcommand{\runsubsection}[1]{{\color{subheadings}\bfseries\large #1}}
\newcommand{\sectionsep}{\vspace{3pt}}

% Optimized Tight Itemize Layout Engine
\newenvironment{tightemize}{
  \begin{itemize}[leftmargin=3mm, topsep=1pt, itemsep=2pt, parsep=0pt]
}{
  \end{itemize}
}

\begin{document}
\color{primarytext}

%======================================================================
% HEADER / TITLE SECTION
%======================================================================
\begin{center}
    {\Huge ALEXANDER} {\Huge\bfseries WRIGHT} \\ \vspace{3pt}
    \color{graytext} \small
    \href{https://alexwright.dev}{alexwright.dev} $|$
    \href{mailto:alex.wright@email.com}{alex.wright@email.com} $|$
    415.555.0192 $|$
    \href{mailto:awright@stanford.edu}{awright@stanford.edu}
\end{center}
\vspace{4pt}

%======================================================================
% TWO COLUMN INTERFACE BLOCK
%======================================================================
\columnratio{0.34}
\setlength{\columnsep}{20pt}
\begin{paracol}{2} % COLUMN ONE: LEFT SIDEBAR

\section{Education}

\runsubsection{Stanford University} \\
\descript{MS in Computer Science} \\
\location{Expected Dec 2027 $|$ Stanford, CA}
\sectionsep

\runsubsection{Oakridge Academy} \\
\location{Grad. May 2022 $|$ Boston, MA}
\sectionsep

\section{Links}
{\small
GitHub:// \href{https://github.com}{\bf alexwright-dev} \\
LinkedIn:// \href{https://linkedin.com}{\bf alex-j-wright} \\
YouTube:// \href{https://youtube.com}{\bf WrightCodes} \\
Twitter:// \href{https://twitter.com}{\bf @alexwright\_cs} \\
Medium:// \href{https://medium.com}{\bf @wright-bytes}
}
\sectionsep

\section{Coursework}
\runsubsection{Graduate} \\
{\small
Deep Learning Architectures \\
Distributed Cloud Storage \\
}
\sectionsep

\runsubsection{Undergraduate} \\
{\small
Analysis of Algorithms \\
Operating Systems Kernels \\
Artificial Intelligence Systems \\
Object-Oriented Programming \\
Database Implementation \\
{\footnotesize \textit{\textbf{(Head Teaching Assistant 3x) }}}
}
\sectionsep

\section{Skills}
\runsubsection{Programming} \\
\location{Over 5000 lines:}
{\small Python $\cdot$ Go $\cdot$ C++ $\cdot$ TypeScript $\cdot$ SQL $\cdot$ \LaTeX} \\
\vspace{2pt}
\location{Over 1000 lines:}
{\small Java $\cdot$ Rust $\cdot$ C $\cdot$ Bash $\cdot$ Dockerfile $\cdot$ HTML} \\
\vspace{2pt}
\location{Familiar:}
{\small Assembly $\cdot$ AWS $\cdot$ PyTorch $\cdot$ Kubernetes}
\sectionsep

\section{Awards}
{\small
2026 -- 1\textsuperscript{st}/120: Stanford Intra-Collegiate Hackathon Grand Champion \\
2025 -- Finalist: Citadel West Coast Datathon Regional Finalist \\
2024 -- National Science Foundation Undergraduate Scholar \\
2023 -- 3\textsuperscript{rd}/80: ACM-ICPC Regional Collegiate Programming Contest \\
2022 -- USAMO Qualifier
}
\sectionsep

\switchcolumn % COLUMN TWO: MAIN BODY CONTENT

\section{Experience}

\runsubsection{Stripe} \descript{$|$ Software Engineering Intern} \\
\location{May 2026 -- Aug 2026 $|$ San Francisco, CA}
\begin{tightemize}
    \item Selected into the ultra-competitive Global Engineering Fellows cohort out of 4,000 global applicants.
    \item Architected and launched Vanguard, an internal tracking framework handling live payment flow validations.
    \item Built reliable high-throughput microservices using Go, processing up to 14,000 transactions per second.
\end{tightemize}
\sectionsep

\runsubsection{Nvidia} \descript{$|$ Accelerated Computing Intern} \\
\location{June 2025 -- Sep 2025 $|$ Santa Clara, CA}
\begin{tightemize}
    \item Programmed core optimizations in the hardware memory layers using C++ and parallelized CUDA layouts.
    \item Designed an experimental caching system reducing latency profiles for large language processing graphs by 14\%.
    \item Collaborated with system layout teams to stress-test high-density database clusters under full load.
\end{tightemize}
\sectionsep

\runsubsection{OpenSource Core} \descript{$|$ Lead Maintainer} \\
\location{Jan 2024 -- May 2025 $|$ Remote Deployment}
\begin{tightemize}
    \item Managed a team of 14 international developers maintaining a distributed asset routing system tool.
    \item Engineered cross-platform containerization scripts reducing onboarding times for fresh environments by 35\%.
\end{tightemize}
\sectionsep

\section{Research}

\runsubsection{Stanford AI Laboratory} \descript{$|$ Head Researcher} \\
\location{Sep 2024 -- Present $|$ Stanford, CA} \\
{\small Co-developed OmniPath alongside research staff, constructing a robust reinforcement learning package that tracks real-time navigation behaviors in dynamic environments. Implemented deep evaluation loops using PyTorch clusters.}
\sectionsep

\runsubsection{Stanford Systems Group} \descript{$|$ Researcher} \\
\location{Jan 2023 -- June 2024 $|$ Stanford, CA} \\
{\small Optimized parallel file system drivers on high-performance compute arrays, cutting average runtime overhead margins down to historic minimum levels across production tests.}
\sectionsep

\section{Publications}
{\small
\noindent Wright, A., et al. (2026). "Predictive Spatial Mapping via Asymmetric Deep Networks." \textit{International Conference on Machine Learning (ICML)}. \\ \vspace{2pt}
\noindent Wright, A., Systems Group. (2025). "Low-Latency File Allocation Mechanisms on Distributed SSD Clusters." \textit{IEEE Systems Journal}.
}
\sectionsep

\end{paracol}
\end{document}
